\documentclass[10pt]{article}
\usepackage[margin=0.85in]{geometry}
\usepackage{amsmath, amssymb}
\usepackage{booktabs}
\usepackage{enumitem}
\usepackage{hyperref}
\usepackage{array}

\newcommand{\bigO}{O}
\newcommand{\SA}{\operatorname{SA}}
\newcommand{\LCP}{\operatorname{LCP}}
\newcommand{\lca}{\operatorname{lca}}
\newcommand{\tin}{\operatorname{tin}}
\newcommand{\tout}{\operatorname{tout}}

\title{EDA Final - Ruta Racso desde Cero\\
\large Solo problemas relevantes del grupo Utec CS3014 Projects}
\author{}
\date{}

\begin{document}
\maketitle

\section*{Que contiene este documento}
Este documento queda limitado a los problemas del grupo de Codeforces que si
sirven para el final segun la lista visible del curso:

\begin{itemize}[leftmargin=*]
\item Proyecto 5
\item Semana 13 - Clase 1 - Offline Dynamic Connectivity
\item Semana 12 - Clase 1 - HLD
\item Semana 11 - Problemas de strings
\item Proyecto 4
\item Semana 10 - Clase 1 - RMQ
\end{itemize}

Se quitaron Proyecto 3/BST implicito y Proyecto 2/Efecto cascada porque no son
prioridad para el final actual. La meta no es practicar todo el curso, sino
aprender desde cero lo que puede aparecer en el examen.

\section*{Orden de estudio si empiezas casi desde cero}
\begin{enumerate}[leftmargin=*]
\item \textbf{RMQ base}: aprende sparse table y por que responde minimo en
$\bigO(1)$.
\item \textbf{LCA/Euler tour}: entiende como convertir arboles en intervalos y
como usar $\lca$ para caminos.
\item \textbf{Strings base}: suffix array, LCP, busqueda de patrones,
subcadenas distintas y longest common substring.
\item \textbf{Tree queries}: DSU on tree, HLD, mascaras de paridad y centroid
decomposition solo a nivel conceptual.
\item \textbf{Offline dynamic connectivity}: segment tree sobre tiempo + DSU
con rollback.
\item \textbf{Proyectos 4 y 5}: usalos como integracion de distancias/caminos
en arboles.
\end{enumerate}

\section*{Mapa limpio de problemas}
\begin{center}
\begin{tabular}{p{0.18\linewidth}p{0.12\linewidth}p{0.38\linewidth}p{0.20\linewidth}}
\toprule
Tema & Contest & Problema & Prioridad \\
\midrule
RMQ & 694840/A & Static RMQ & Maxima \\
Strings & 695758/A & Suffix array y LCP array & Maxima \\
Strings & 695758/B & Busqueda de subcadenas & Maxima \\
Strings & 695758/C & Contando subcadenas & Alta \\
Strings & 695758/D & Subcadenas diferentes & Maxima \\
Strings & 695758/E & Subcadena en comun mas larga & Maxima \\
Strings & 695758/F & Ordenando subcadenas & Alta \\
Strings & 695758/G & Bordes & Alta \\
Strings & 695758/H & Refran & Alta \\
Strings & 695758/I & Substring periodica & Media/alta \\
Arboles/HLD & 697527/A & Tree and Queries & Maxima \\
Arboles/HLD & 697527/B & Arpa letter-marked tree & Alta \\
Arboles/HLD & 697527/C & Laminar Family & Media \\
Arboles/HLD & 697527/D & Xenia and Tree & Alta \\
Arboles/HLD & 697527/E & Beauty and The Tree & Media \\
Dynamic graph & 698987/A & Connect and Disconnect & Maxima \\
Offline intervals & 698987/B & Addition on Segments & Alta \\
Proyecto final & 699336/A & Visitas & Alta \\
Distancias & 694842/A & Distancia unitaria & Media \\
\bottomrule
\end{tabular}
\end{center}

\newpage
\section{Semana 10 - Static RMQ}

\subsection*{Problema A: Static RMQ}
Link:
\url{https://codeforces.com/group/AmklRC8lhZ/contest/694840/problem/A}

\subsection*{Que debes aprender primero}
RMQ significa \emph{Range Minimum Query}. Te dan un arreglo fijo $A$ y muchas
consultas $[l,r]$. Debes responder:
\[
\min(A[l], A[l+1], \ldots, A[r]).
\]
La palabra clave es \textbf{static}: el arreglo no cambia. Si el arreglo cambia,
necesitarias segment tree o Fenwick segun la operacion. Como no cambia, puedes
preprocesar fuerte.

\subsection*{Solucion principal: sparse table}
Define una tabla:
\[
st[k][i] = \min(A[i], A[i+1], \ldots, A[i+2^k-1]).
\]
Es decir, guarda el minimo de cada bloque de longitud potencia de dos.

La transicion es:
\[
st[k][i] = \min(st[k-1][i], st[k-1][i+2^{k-1}]).
\]
El bloque grande de longitud $2^k$ se parte en dos bloques de longitud
$2^{k-1}$.

Para responder $[l,r]$, toma:
\[
k=\lfloor \log_2(r-l+1)\rfloor.
\]
Entonces dos bloques de longitud $2^k$ cubren todo el intervalo:
\[
[l,l+2^k-1] \quad \text{y} \quad [r-2^k+1,r].
\]
La respuesta es:
\[
\min(st[k][l], st[k][r-2^k+1]).
\]

\subsection*{Por que funciona}
Los dos bloques pueden solaparse, pero eso no importa porque la operacion
$\min$ es idempotente:
\[
\min(x,x)=x.
\]
Mientras ambos bloques cubran todo el rango, el minimo global esta en alguno de
ellos. Si el minimo cae en la zona solapada, contarlo dos veces no cambia nada.

\subsection*{Complejidad}
Preprocesamiento:
\[
\bigO(n\log n).
\]
Consulta:
\[
\bigO(1).
\]
Espacio:
\[
\bigO(n\log n).
\]
Para teoria avanzada, 6.851 muestra RMQ en $\bigO(n)$ espacio y $\bigO(1)$
query usando bloques macro/micro, pero para resolver y explicar en examen,
sparse table es la base que debes dominar.

\newpage
\section{Semana 11 - Strings}

\subsection*{Base comun: suffix array y LCP}
Un sufijo de $S$ es una subcadena que empieza en una posicion $i$ y termina al
final. El suffix array $\SA$ es el arreglo de posiciones de sufijos ordenados
lexicograficamente. Si:
\[
\SA = [p_0,p_1,\ldots,p_{n-1}],
\]
entonces:
\[
S[p_0..] < S[p_1..] < \cdots < S[p_{n-1}..].
\]

El arreglo $\LCP$ guarda cuanto prefijo comun tienen sufijos consecutivos:
\[
\LCP[i] = \operatorname{lcp}(S[\SA[i-1]..], S[\SA[i]..]).
\]

Esto sirve porque muchos problemas de strings se vuelven problemas de ordenar
sufijos y medir prefijos comunes.

\subsection{A. Suffix array y LCP array}
Link:
\url{https://codeforces.com/group/AmklRC8lhZ/contest/695758/problem/A}

\paragraph{Que saber.}
Debes poder construir y explicar $\SA$ y $\LCP$. El algoritmo de clase mas
importante es Manber--Myers:

\begin{enumerate}[leftmargin=*]
\item Primero ordenas sufijos por su primer caracter.
\item Luego por los primeros $2$ caracteres.
\item Luego por los primeros $4$.
\item Luego por $8$, $16$, etc.
\end{enumerate}

En cada paso, el sufijo que empieza en $i$ se representa por:
\[
(rank[i], rank[i+2^k]).
\]
Si ordenas esos pares, ordenas por prefijos de longitud $2^{k+1}$.

\paragraph{Correctitud.}
Si ya sabes el orden relativo de todas las subcadenas de longitud $2^k$, entonces
una subcadena de longitud $2^{k+1}$ queda determinada por sus dos mitades:
primera mitad y segunda mitad. Por eso ordenar pares de ranks es suficiente.

\paragraph{Complejidad.}
Hay $\bigO(\log n)$ rondas y cada ronda puede hacerse en $\bigO(n)$ con counting
sort/radix sort sobre ranks. Total:
\[
\bigO(n\log n).
\]

\subsection{B. Busqueda de subcadenas}
Link:
\url{https://codeforces.com/group/AmklRC8lhZ/contest/695758/problem/B}

\paragraph{Idea.}
Para buscar un patron $P$ dentro de un texto fijo $S$, construyes el suffix
array de $S$. Todas las ocurrencias de $P$ son sufijos que empiezan con $P$.
En orden lexicografico, todos esos sufijos aparecen consecutivos.

\paragraph{Como se responde.}
Haces dos busquedas binarias:
\begin{itemize}[leftmargin=*]
\item primera posicion del suffix array cuyo sufijo es $\ge P$;
\item primera posicion cuyo sufijo es $> P$ como prefijo.
\end{itemize}
El intervalo entre ambas posiciones contiene todas las ocurrencias.

\paragraph{Por que no usar KMP siempre.}
KMP es excelente para un patron. Si tienes muchas consultas sobre el mismo texto,
suffix array preprocesa el texto una vez y luego cada patron se busca rapido.

\subsection{C. Contando subcadenas}
Link:
\url{https://codeforces.com/group/AmklRC8lhZ/contest/695758/problem/C}

Este problema entrena la idea de que las ocurrencias de un patron se cuentan por
el tamano del intervalo de suffix array donde los sufijos tienen ese patron como
prefijo. Si el intervalo es $[L,R]$, la cantidad de ocurrencias es:
\[
R-L+1.
\]

Si la variante pide contar subcadenas con alguna propiedad, piensa primero:
\begin{itemize}[leftmargin=*]
\item si depende de apariciones de un patron, usa intervalo en $\SA$;
\item si depende de subcadenas distintas, usa suma de LCP;
\item si depende de repeticion, usa minimos en rangos de LCP.
\end{itemize}

\subsection{D. Subcadenas diferentes}
Link:
\url{https://codeforces.com/group/AmklRC8lhZ/contest/695758/problem/D}

El numero total de subcadenas de un string de longitud $n$ es:
\[
\frac{n(n+1)}2.
\]
Pero hay repetidas. Al ordenar sufijos, cada sufijo aporta sus prefijos. El
sufijo $\SA[i]$ aporta $n-\SA[i]$ prefijos, pero los primeros $\LCP[i]$ ya
aparecieron en el sufijo anterior.

Formula:
\[
\#\text{subcadenas distintas}
= \frac{n(n+1)}2 - \sum_i \LCP[i].
\]

\paragraph{Ejemplo mental.}
Si dos sufijos vecinos comparten prefijo de longitud $3$, entonces las subcadenas
de longitud $1$, $2$ y $3$ que empiezan en el segundo sufijo ya fueron contadas
por el primero.

\subsection{E. Subcadena en comun mas larga}
Link:
\url{https://codeforces.com/group/AmklRC8lhZ/contest/695758/problem/E}

Dado $A$ y $B$, construye:
\[
T = A \# B \$,
\]
donde $\#$ y $\$$ no aparecen en los strings. Construyes $\SA$ y $\LCP$ de $T$.
La respuesta es el maximo $\LCP[i]$ tal que los sufijos $\SA[i-1]$ y $\SA[i]$
vienen de strings distintos.

\paragraph{Por que basta revisar vecinos.}
Si una subcadena aparece en ambos strings, todos los sufijos que empiezan con esa
subcadena quedan juntos en el suffix array. Dentro de ese bloque, si hay sufijos
de ambos strings, existe algun par vecino de origen distinto.

\subsection{F. Ordenando subcadenas}
Link:
\url{https://codeforces.com/group/AmklRC8lhZ/contest/695758/problem/F}

Para comparar dos subcadenas $S[l_1,r_1]$ y $S[l_2,r_2]$, necesitas saber su
primer caracter distinto. Eso se obtiene con:
\[
\operatorname{lcp}(l_1,l_2).
\]
Para obtener LCP entre sufijos arbitrarios, usas RMQ sobre el arreglo $\LCP$:
si los sufijos estan en posiciones $a$ y $b$ del suffix array, entonces:
\[
\operatorname{lcp}(a,b)=\min(\LCP[a+1],\ldots,\LCP[b]).
\]

\paragraph{Conclusion.}
Ordenar subcadenas requiere combinar strings con RMQ. Esta es una conexion muy
probable para examen: suffix array + LCP + RMQ.

\subsection{G. Bordes}
Link:
\url{https://codeforces.com/group/AmklRC8lhZ/contest/695758/problem/G}

Un borde es una cadena que es prefijo y sufijo a la vez. Para estudiar desde
cero, usa KMP:
\[
\pi[i] = \text{longitud del mayor borde de } S[0..i].
\]
Los bordes de todo $S$ se obtienen siguiendo:
\[
k=\pi[n-1], \quad k=\pi[k-1], \quad \ldots
\]
hasta llegar a cero.

\paragraph{Relacion con suffix array.}
Tambien puedes verlo como LCP entre el sufijo que empieza en $0$ y otros sufijos
que terminan justo al final, pero KMP es mas simple para bordes.

\subsection{H. Refran}
Link:
\url{https://codeforces.com/group/AmklRC8lhZ/contest/695758/problem/H}

Este tipo de problema suele buscar una subcadena repetida que maximiza una
funcion como:
\[
\text{longitud} \times \text{frecuencia}.
\]
En suffix array, una subcadena aparece muchas veces si hay muchos sufijos
consecutivos que comparten un prefijo largo. El arreglo $\LCP$ se interpreta
como un histograma: un minimo alto en un bloque grande significa una subcadena
larga que aparece muchas veces.

\paragraph{Estructura mental.}
Usa pila monotona sobre $\LCP$ para encontrar, para cada altura $h$, el bloque
mas grande donde todos los LCP son al menos $h$. Eso da frecuencia y longitud.

\subsection{I. Substring periodica}
Link:
\url{https://codeforces.com/group/AmklRC8lhZ/contest/695758/problem/I}

Una cadena tiene periodo $p$ si:
\[
S[i]=S[i+p]
\]
siempre que ambas posiciones existan. Para substrings, necesitas verificar
igualdades largas entre posiciones. Herramientas:
\begin{itemize}[leftmargin=*]
\item Z-function o prefix-function para periodos globales;
\item LCP entre sufijos para comparar substrings rapidamente;
\item RMQ sobre LCP para LCP arbitrario.
\end{itemize}

Este es el mas avanzado del bloque. Si no sabes nada, primero asegura A--E.

\newpage
\section{Semana 12 - HLD y Tree Queries}

\subsection*{Base: caminos y subarboles}
En un arbol, hay dos tipos de preguntas muy comunes:
\begin{itemize}[leftmargin=*]
\item sobre un subarbol;
\item sobre el camino entre dos nodos.
\end{itemize}

Para subarboles, Euler tour convierte el subarbol de $u$ en el intervalo
$[\tin(u),\tout(u)]$. Para caminos, usas LCA/HLD.

\subsection{A. Tree and Queries}
Link:
\url{https://codeforces.com/group/AmklRC8lhZ/contest/697527/problem/A}

Problema: cada nodo tiene color; una consulta pregunta en el subarbol de $v$
cuantos colores aparecen al menos $k$ veces.

\paragraph{Solucion.}
Usa DSU on tree. Conservas la informacion del hijo pesado y agregas los hijos
ligeros solo cuando hace falta. Mantienes:
\[
cntColor[c] = \text{apariciones del color }c,
\]
y una estructura para saber cuantos colores tienen frecuencia al menos $k$.

\paragraph{Por que funciona.}
Antes de responder consultas de $v$, la bolsa contiene exactamente su subarbol.
Los hijos son disjuntos; si agregas todos los hijos y el nodo, tienes el subarbol
completo.

\paragraph{Complejidad.}
Cada nodo entra a bolsas asociadas a aristas ligeras $\bigO(\log n)$ veces. La
cota segura es $\bigO(n\log n)$.

\subsection{B. Arpa's letter-marked tree and Mehrdad's Dokhtar-kosh paths}
Link:
\url{https://codeforces.com/group/AmklRC8lhZ/contest/697527/problem/B}

La idea central es palindromos por paridad. Una multiconjunto de letras puede
formar un palindromo si como maximo una letra aparece cantidad impar de veces.
Representa paridades con una mascara de bits.

Para un camino, la mascara se obtiene con xor de prefijos desde la raiz. Dos
mascaras son compatibles si:
\begin{itemize}[leftmargin=*]
\item son iguales; o
\item difieren en exactamente un bit.
\end{itemize}

Esto suele combinarse con DSU on tree para contar pares dentro de subarboles.

\subsection{C. Laminar Family}
Link:
\url{https://codeforces.com/group/AmklRC8lhZ/contest/697527/problem/C}

Una familia laminar cumple que dos conjuntos cualesquiera son disjuntos o uno
contiene al otro. En arboles, los subarboles tienen esta propiedad cuando se
representan como intervalos Euler:
\[
subtree(u) = [\tin(u),\tout(u)].
\]
Dos intervalos de subarbol no se cruzan parcialmente: o son disjuntos o uno
contiene al otro. La teoria importante es saber probar inclusion con tiempos:
\[
u \text{ ancestro de } v
\iff \tin(u)\le \tin(v)\le \tout(u).
\]

\subsection{D. Xenia and Tree}
Link:
\url{https://codeforces.com/group/AmklRC8lhZ/contest/697527/problem/D}

Problema clasico: hay nodos rojos y consultas por distancia al rojo mas cercano.
La estructura esperada es centroid decomposition.

\paragraph{Desde cero.}
Un centroide parte el arbol en componentes de tamano a lo mucho la mitad. Si
repites recursivamente, cada nodo tiene $\bigO(\log n)$ centroides ancestros.

\paragraph{Operacion pintar.}
Para un nodo rojo $v$, actualizas cada centroide ancestro $c$ con:
\[
best[c] = \min(best[c], dist(v,c)).
\]

\paragraph{Operacion preguntar.}
Para un nodo $u$, pruebas todos sus centroides ancestros:
\[
\min_c best[c] + dist(u,c).
\]
Como hay $\bigO(\log n)$ centroides ancestros, cada operacion cuesta
$\bigO(\log n)$ si las distancias se obtienen con LCA o precomputacion.

\subsection{E. Beauty and The Tree}
Link:
\url{https://codeforces.com/group/AmklRC8lhZ/contest/697527/problem/E}

Este problema es menos central que A--D. Entra como arbol + DP/probabilidad.
Si aparece algo parecido en examen, la ruta mental es:
\begin{enumerate}[leftmargin=*]
\item enraizar el arbol;
\item calcular tamanos de subarbol;
\item definir que informacion sube de hijos a padre;
\item si la respuesta depende de ``ver el arbol desde todos los nodos'', usar
rerooting.
\end{enumerate}

\newpage
\section{Semana 13 - Offline Dynamic Connectivity}

\subsection{A. Connect and Disconnect}
Link:
\url{https://codeforces.com/group/AmklRC8lhZ/contest/698987/problem/A}

Este es core para dynamic graph. Hay operaciones que conectan y desconectan
aristas, y consultas de conectividad. DSU normal no puede borrar aristas, por
eso se procesa offline.

\paragraph{Paso 1: convertir aristas en intervalos de vida.}
Si una arista se agrega en tiempo $l$ y se elimina en tiempo $r$, entonces vive
en el intervalo $[l,r)$. Si nunca se elimina, vive hasta el final.

\paragraph{Paso 2: segment tree sobre tiempo.}
Cada intervalo de vida se descompone en $\bigO(\log q)$ nodos de un segment tree
que cubre los tiempos.

\paragraph{Paso 3: DFS con DSU rollback.}
Al entrar a un nodo del segment tree, agregas todas sus aristas al DSU. Al salir,
deshaces esas uniones. En una hoja, el DSU contiene exactamente las aristas
activas en ese instante.

\paragraph{Por que rollback.}
Path compression hace dificil deshacer cambios. Por eso DSU rollback usa union
by size/rank, guarda cambios en una pila, y revierte al salir del intervalo.

\paragraph{Complejidad.}
Cada arista aparece en $\bigO(\log q)$ nodos. Cada union/find cuesta
$\bigO(\log n)$ sin compresion o casi constante en la practica con union by size.
La cota segura:
\[
\bigO((q+m)\log q \log n).
\]

\subsection{B. Addition on Segments}
Link:
\url{https://codeforces.com/group/AmklRC8lhZ/contest/698987/problem/B}

Aunque no sea grafo, entrena la misma idea: una operacion existe en un intervalo.
Cuando veas ``agregar algo en un segmento de tiempo'' o ``activar/desactivar'',
piensa en segment tree sobre tiempo.

La leccion importante para el examen:
\begin{itemize}[leftmargin=*]
\item si hay operaciones que se activan y desactivan;
\item y puedes leer todo el input antes de responder;
\item entonces convierte cada objeto en intervalos de vida.
\end{itemize}

\newpage
\section{Proyecto 5 y Proyecto 4}

\subsection{Proyecto 5: Visitas}
Link:
\url{https://codeforces.com/group/AmklRC8lhZ/contest/699336/problem/A}

Por estar al final del bloque y por el nombre, debe tratarse como problema de
arboles/caminos/visitas. Si es una variante tipo consultas sobre caminos, las
herramientas probables son:
\begin{itemize}[leftmargin=*]
\item LCA para descomponer caminos;
\item diferencias en arbol si las visitas son offline;
\item HLD si hay updates/queries sobre caminos;
\item DSU on tree si las visitas son por subarbol/color.
\end{itemize}

Como no debes aprender todo, la regla practica es:
\begin{enumerate}[leftmargin=*]
\item Si preguntan cantidad de rutas que pasan por aristas: diferencia en arbol.
\item Si preguntan max/min/suma en camino con cambios: HLD.
\item Si preguntan colores en subarbol: DSU on tree.
\item Si preguntan conectividad con cambios: offline dynamic connectivity.
\end{enumerate}

\subsection{Proyecto 4: Distancia unitaria}
Link:
\url{https://codeforces.com/group/AmklRC8lhZ/contest/694842/problem/A}

Si es de distancias en arbol, la formula que debes memorizar es:
\[
dist(u,v)=depth[u]+depth[v]-2depth[\lca(u,v)].
\]
Si es de grafo no ponderado, la palabra ``unitaria'' sugiere BFS porque todas
las aristas pesan 1. Para el final, este problema vale como practica secundaria:
distancia en arbol/grafo, no como tema principal.

\newpage
\section{Resumen final: que estudiar y que ignorar}

\subsection*{Estudia}
\begin{itemize}[leftmargin=*]
\item Static RMQ.
\item Suffix array, LCP, busqueda de patrones, subcadenas distintas, LCS,
bordes y periodicidad.
\item DSU on tree, Euler tour, LCA, mascaras xor en arboles.
\item Centroid decomposition para Xenia and Tree a nivel de idea.
\item Offline dynamic connectivity con segment tree sobre tiempo y DSU rollback.
\item Proyecto 5 y Proyecto 4 solo como integracion de arboles/distancias.
\end{itemize}

\subsection*{Ignora para esta ruta}
\begin{itemize}[leftmargin=*]
\item Proyecto 3 - BST implicito.
\item Proyecto 2 - Efecto cascada.
\item Cualquier problema externo que no este en esta lista, salvo que el profesor
lo mencione directamente.
\end{itemize}

\section*{Checklist para resolver en papel}
\begin{enumerate}[leftmargin=*]
\item Identifica si es string, arbol estatico, camino en arbol, subarbol, o
grafo dinamico.
\item Escribe la estructura: sparse table, suffix array, DSU on tree, HLD,
centroid decomposition, o DSU rollback.
\item Escribe el invariante: que guarda la estructura y por que representa el
problema.
\item Prueba correctitud en 5-8 lineas.
\item Cierra con complejidad de preprocesamiento, query, update y memoria.
\end{enumerate}

\end{document}
